\documentclass[letterpaper,twocolumn,10pt]{article}
\usepackage{usenix2019_v3}

% to be able to draw some self-contained figs
\usepackage{tikz}
\usepackage{amsmath}

% inlined bib file
\usepackage{filecontents}

%-------------------------------------------------------------------------------
\begin{document}
%-------------------------------------------------------------------------------

%don't want date printed
\date{}

% make title bold and 14 pt font (Latex default is non-bold, 16 pt)
\title{\Large \bf Literature Review}

\author{
  {\rm Aayuman Ghangurde}\\
  Purdue University
  \and
  {\rm Sai Palaparthi}\\
  Purdue University
  \and
  {\rm Rakib Ahsan}\\
  Purdue University
  \and
  {\rm Zeke Ulrich}\\
  Purdue University
}

\maketitle

%-------------------------------------------------------------------------------
\section{Related Work}
%-------------------------------------------------------------------------------


\subsection{Technical Report on a Virtual CTAP2 WebAuthn Authenticator~\cite{culnane2021virtualctap2webauthn}}

\paragraph{Summary:}
This technical report covers the development of the associated open source
Virtual CTAP2 WebAuthn authenticator. In addition to
a software only authenticator it also provides a proof of concept implementation
of a Trusted Platform Module (TPM) based authenticator, with associated
interfaces and libraries for using a TPM as the underlying credential store.

\paragraph{Shortcomings:}
This project wants several improvements we hope to make, including:
\begin{itemize}
  \item the authors deliberately keep client-PIN verification outside
        the TPM (software ES256 provider), reasoning that the PIN is typed
        anyway so it "doesn't warrant" TPM protection. This is a waste of
        modern TPMs' built-in dictionary-attack prevention, which we can enable.
  \item their authenticator uses a flat key hierarchy with no dictionary-attack
        protection on the intermediate key. If we use passkey signing keys
        as children of an intermediate PIN key that has a lockout policy,
        then unlocking the PIN key makes them usable just as with Windows
        Hello and we have brute force protection against the PIN.
  \item the paper notes that their method of emulating a USB HID device
        is fragile. We can build the authenticator as a PAM module for
        extensibility.
  \item the authenticator suuports only ES256, which we could expand to
        RSA and Ed25519.
\end{itemize}

\subsection{Securing FIDO2 Credentials Using Intel SGX~\cite{matos2026securing}}

\paragraph{Summary:}
This paper builds a centralized, replicated FIDO2 
credential manager similar to the one we're building, 
using /dev/uhid to emulate a USB HID FIDO2 device like
tpm-fido/VirtualWebAuthn does.

\paragraph{Shortcomings:}
Their conclusion states they want to 
"improve the Authenticator Emulator to use a local vault 
(e.g., TPM-based) to store a UID protected by a secret
or biometrics." In fact, Intel has discontinued SGX
so the credential manager must be migrated to a TPM anyway. 

\subsection{The State of Passkeys: Studying the Adoption and Security of Passkeys on the Web~\cite{jannett2026stateofpasskeys}}

\paragraph{Summary:} The authors perform a comprehensive
security evaluation of passkey implementation in websites.
The authors identify lack of standardization and serious
security issues in a majority of websites, including use
of deprecated cryptographic algorithms and critical or
high CVSS scores. The authors present
\texttt{PASSKEYS-RADAR} for tracking deployment of
passkeys on the web, and \texttt{PASSKEYS-ATTACKER} for
exploting security vulnerabilities on insecure sites.

\paragraph{Shortcomings:} The authors acknowledge
several limitations of their approach, including:
\begin{itemize}
  \item \emph{Overestimation of adoption rate}: the automated detec-
        tion of passkey-enabled websites inevitably produces false
        positives, which leads to an overestimation of the adoption.
  \item \emph{Limited representativeness of the dataset}: the
        dataset is is biased toward publicly accessible, consumer-
        facing services. Certain sectors, such as online banking, could
        not be systematically evaluated due to the lack of suitable test
        accounts.
  \item \emph{Different authentication backends}: It is not possible to
        automatically determine whether multiple domains operated
        by the same provider (e.g., gmail.com and youtube.com)
        share a common authentication backend.
\end{itemize}

\subsection{FIDO2, CTAP 2.1, and WebAuthn 2: Provable Security and Post-Quantum Instantiation~\cite{bindel2023fido2ctap21webauthn2}}

\paragraph{Summary:} The authors formally
analyze FIDO2's security with the
CTAP 2.1 and WebAuthn 2 sub-protocols.
The authors prove post-quantum security for
FIDO2 under certain conditions and minor protocol extensions,
and show a simple improvement for the downgrade
resilience of FIDO2.

\paragraph{Shortcomings:}
Since the paper is an abstract proof of
CTAP and WebAuthn, the authors
do not model some of the new CTAP 2.1 features for enterprise
customers, and do not make formal statements about the
unlinkability of credentials or other detailed privacy statements.
Thus, the implementation of these standards may have
unaddressed vulnerabilities.



%-------------------------------------------------------------------------------
\bibliographystyle{plain}
\bibliography{\jobname}

%%%%%%%%%%%%%%%%%%%%%%%%%%%%%%%%%%%%%%%%%%%%%%%%%%%%%%%%%%%%%%%%%%%%%%%%%%%%%%%%
\end{document}
%%%%%%%%%%%%%%%%%%%%%%%%%%%%%%%%%%%%%%%%%%%%%%%%%%%%%%%%%%%%%%%%%%%%%%%%%%%%%%%%

%%  LocalWords:  endnotes includegraphics fread ptr nobj noindent
%%  LocalWords:  pdflatex acks